\documentclass[10pt]{article}
\usepackage[T1]{fontenc}
\usepackage[utf8]{inputenc}
\usepackage[french]{babel}
\usepackage{newpxtext,newpxmath}
\usepackage[a4paper,top=1.5cm,bottom=1.35cm,left=1.65cm,right=1.65cm]{geometry}
\usepackage{microtype}
\usepackage{xcolor}
\usepackage{titlesec}
\usepackage[most]{tcolorbox}
\usepackage{enumitem}
\usepackage{array}
\usepackage{booktabs}
\usepackage{fancyhdr}

\definecolor{navy}{HTML}{1F3A5F}
\definecolor{grey}{HTML}{54606E}
\definecolor{green}{HTML}{2C6E49}
\definecolor{amber}{HTML}{9C6A16}
\definecolor{red}{HTML}{A6322E}
\definecolor{plum}{HTML}{7A4EA8}
\definecolor{boxgray}{HTML}{F3F5F8}
\definecolor{boxred}{HTML}{FCF2F1}
\definecolor{boxedge}{HTML}{D9E0EA}

\setlength{\parindent}{0pt}
\setlength{\parskip}{2.6pt}
\linespread{1.0}

\pagestyle{fancy}\fancyhf{}
\lhead{\small\color{grey}Scella}
\rhead{\small\color{grey}Composants -- ce qu'il reste à faire}
\cfoot{\small\color{grey}\thepage}
\renewcommand{\headrulewidth}{0.3pt}

\titleformat{\section}{\color{navy}\large\bfseries}{}{0pt}{}
\titlespacing{\section}{0pt}{9pt}{3pt}
\titleformat{\subsection}{\color{navy}\bfseries}{}{0pt}{}
\titlespacing{\subsection}{0pt}{7pt}{1.5pt}

\newtcolorbox{these}{colback=boxgray,colframe=boxedge,boxrule=0.5pt,arc=2.5pt,
  left=8pt,right=8pt,top=4.5pt,bottom=4.5pt}
\newtcolorbox{alerte}{colback=boxred,colframe=red!55,boxrule=0.8pt,arc=2.5pt,
  left=8pt,right=8pt,top=4.5pt,bottom=4.5pt}

\newcommand{\ok}{{\color{green}\textbf{SERVI}}}
\newcommand{\partiel}{{\color{amber}\textbf{PARTIEL}}}
\newcommand{\design}{{\color{grey}\textbf{CONÇU}}}
\newcommand{\nope}{{\color{red}\textbf{NON SERVI}}}
\newcommand{\sz}[1]{{\small\color{grey}\textsf{#1}}}
\newcommand{\code}[1]{\texttt{\small #1}}

\begin{document}

\begin{center}
{\Huge\bfseries\color{navy} Scella}\\[4pt]
{\large\color{grey}\itshape Analyse complète des composants : ce qu'il reste à faire pour chacun}\\[4pt]
{\small\color{grey} 5 août 2026 · v2 \; · \; audit de code, 6 cartographes en parallèle \; · \; 1,84 M tokens de lecture, 549 appels d'outils}
\end{center}
\vspace{2pt}

\begin{these}
\textbf{Le résultat central.} \emph{Presque tout est construit ; presque rien n'est servi.} Et le seul
actif réellement différenciant -- \textbf{le reçu} -- porte un défaut d'une ligne qui lui fait renvoyer
\code{verified: true} sur une chaîne tronquée. Le correctif existe déjà, dans le fichier d'à côté.
\end{these}

\section{1 · Méthode et portée}

Cette analyse ne repose sur aucun document de conception pris pour argent comptant. Chaque composant a
été relu \emph{dans le code d'aujourd'hui} par un cartographe dédié, avec pour consigne explicite de
séparer trois natures de reste-à-faire : les \textbf{jalons déclarés} (nommés dans un document, avec
leur libellé exact), les \textbf{trous découverts} qu'aucun document ne nomme, et les \textbf{résidus
assumés} que le code lui-même documente comme connus et tolérés. Les tailles sont notées XS à XL.

Trois de mes affirmations antérieures ont été \textbf{réfutées} par cette relecture ; elles sont
signalées à leur place plutôt que corrigées en silence. C'est la discipline habituelle du projet :
mesurer plutôt que supposer, et dire quand la mesure contredit.

\begin{alerte}
\textbf{Le défaut qui domine tout.} \code{GET /v1/swarm/verify/:id} recalcule la racine de la chaîne et
\textbf{ne la compare jamais} à la racine scellée (\code{routes/verify.rs:41}), alors que la comparaison
est faite correctement, un fichier plus loin, dans le replay R0 (\code{routes/replay.rs:79}).
Conséquence directe : un tiers disposant d'un accès en écriture à la base de données, \textbf{et
d'aucune clé}, peut supprimer les derniers enregistrements d'un run et obtenir malgré tout
\code{verified: true}. C'est \emph{exactement} le modèle de menace contre lequel le reçu est censé
protéger. Le correctif tient en une ligne : ajouter \code{root\_matches} à la réponse et le combiner
au verdict. \sz{Taille S.}
\end{alerte}

\section{2 · Tableau synoptique}

\renewcommand{\arraystretch}{1.12}
\begin{center}\small
\begin{tabular}{@{}rl l >{\raggedright\arraybackslash}p{8.1cm}c@{}}
\toprule
& \textbf{Composant} & \textbf{Statut réel} & \textbf{Ce qui reste} & \textbf{Taille} \\
\midrule
1 & Le Relais (Fil B) & \ok & servi opt-in ; durcir le point d'extension + 4 résidus & \sz{S} \\
2 & \textbf{Le Sandwich} & \design & 0 code, \emph{substrat prêt} : les 4 jalons C.1.a$\rightarrow$d & \sz{L} \\
3 & L'essaim médiateur & \design & 0 code : les 5 jalons E.0$\rightarrow$E.4 & \sz{L} \\
4 & Fil C, la Colonne & \design & 0 code : les 5 jalons C.0$\rightarrow$C.4 & \sz{XL} \\
5 & \textbf{Les nano-robots} & \partiel & priors vides (\code{sin(i)}) ; état initial hors boîte, population ignorée & \sz{S} \\
6 & \textbf{L'essaim $\Phi$} & \partiel & concurrence non pure (non commité) ; 3 dettes de mesure & \sz{M} \\
7 & \textbf{Le méta-essaim} & \nope & \emph{mock par défaut} ; moteur jamais vérifié ; M0 + M2 + M3 + M4 & \sz{XL} \\
8 & Le reçu, pile A & \partiel & non résistant à la troncature ; 13 trous, dont 3 critiques & \sz{M} \\
9 & Le reçu, pile B & \ok & S0$\rightarrow$S6, 245 tests ; \textbf{aucun producteur}, aucun ancrage & \sz{L} \\
10 & Couche 2 servie & \partiel & off par défaut ; PII fail-open, budget no-op, e2e 1/5 & \sz{M} \\
11 & Runtime nano & \nope & rien ne l'exécute ; 7 défauts confirmés + 2 nouveaux graves & \sz{L} \\
12 & Banc de mesure & \partiel & statistiques solides, non régressible ; cassettes record-replay & \sz{L} \\
13 & Flexible MLP & \design & hors dépôt ; validation gold, matériel & \sz{XL} \\
\bottomrule
\end{tabular}
\end{center}

\section{3 · Composant par composant}

\subsection{3.1 · Le Relais (Fil B) -- le seul composant vraiment fini}

C'est la brique aboutie, et elle mérite d'être décrite comme telle sans emphase : les jalons B.0 à B.4
sont complets et \textbf{réellement servis} en opt-in via \code{EFE\_SLM\_REFORMULATOR\_URL}. La
couverture est de 49 tests dédiés, dont 44 tournent hors ligne contre un faux serveur llama monté sur
un \code{TcpListener} réel -- ce ne sont donc pas de simples doublures en mémoire -- et 5 tests live
marqués \code{\#[ignore]}. Quatre revues adverses ont été passées. Le trait asynchrone, les estampilles
de moteur, le disjoncteur à sonde unique, les deadlines de mission, la capture d'audit par hop et le
scellement individuel de chaque hop dans la chaîne : tout cela existe et tourne.

Ce qui reste y est de nature défensive plutôt que constructive. Le point le plus important est que
\emph{le point d'extension est piégé} : la vue qui calcule les compteurs de moteurs comporte un bras
\code{\_ => \{\}} (\code{routes/llm.rs:159}) qui avalerait silencieusement toute variante future, et
l'invariant d'honnêteté \code{slm\_mediated} n'est asserté \textbf{qu'au cas négatif}. Autrement dit :
le garde-fou central du Fil B tient par construction, pas par test. Rendre ce filtrage exhaustif et
écrire les tests de vue manquants coûte une demi-journée et supprime la classe de bug la plus probable
de tout le chantier à venir. Trois résidus documentés subsistent par ailleurs : une double défaillance
(le petit modèle \emph{et} son repli) fait sortir la tentative hors du reçu ; l'URL du petit modèle
n'est pas contrainte au local, si bien que son egress n'est ni gardé ni décompté ; et la concurrence
du disjoncteur n'est assertée que séquentiellement.

\subsection{3.2 · Le Sandwich (l'Encadrement) -- zéro ligne, mais le substrat est prêt}

C'est le composant dont le reste-à-faire est le mieux borné, parce que \emph{tout ce qu'il consommera
existe déjà}. Le compresseur est un client réel de 914 lignes, avec grammaire dynamique et découpage
map-reduce au-delà d'environ 512 tokens estimés ; le vérifieur rend un verdict fermé à quatre valeurs.
Les emplacements \code{compressor} et \code{verifier} sont \textbf{présents} sur le porteur de clients
(\code{coordinator.rs:43,45}), à \code{None} par défaut, et les circuits \code{"compressor"} et
\code{"verifier"} sont \textbf{déjà créés} par \code{default\_slm\_circuits()} au seuil 3 / reset 10 s.
Le protocole d'usage est lui-même éprouvé ailleurs, dans les missions A/B : seuil métier à 250 tokens,
cible proportionnelle, détection de sortie dégénérée, ré-injection des questions, et un vérifieur qui
\emph{annote sans jamais bloquer}. Conséquence directe et importante : \textbf{il n'y a pas de « B.1
bis »} -- aucun asset de grammaire à écrire, aucun modèle à entraîner.

Ce qui reste tient en quatre tranches. \textbf{C.1.a}, le socle, est la seule qui touche le fil scellé :
généralisation additive du type de hop (\code{Prep}, \code{Translate} par défaut, \code{Check}),
extraction du runtime de hop du Fil B vers un module partagé, et surtout une capture qui \emph{compte}
les appels -- car le map-reduce du compresseur casse l'hypothèse implicite « un hop égale une entrée
d'audit ». S'y ajoutent les exports anti-dérive depuis le coordinateur : une seule source de vérité,
jamais une copie, la leçon du garde de souveraineté dupliqué. \sz{S} \textbf{C.1.b} livre le prep servi,
avec l'ordre des gardes inchangé et la validation des sorties dégénérées. \sz{M} \textbf{C.1.c} livre le
check, avec des clamps \emph{scellés} : on ne vérifie jamais en silence un claim tronqué. \sz{S--M}
\textbf{C.1.d} livre la surface et la revue adverse. \sz{S}

Le point de vigilance est mesurable et connu : le timeout interne du compresseur est de \textbf{5
secondes par appel}, et le map-reduce en fait plusieurs -- c'est la deadline de mission qui borne le hop
entier. À mesurer avant d'ajouter le moindre réglage.

\subsection{3.3 · L'essaim médiateur et Fil C -- conçus, rien codé}

Reconfirmé par recherche directe : \code{MediationRegistry}, \code{EdgePlan}, \code{SwarmMediator} et
\code{CanonicalSwarmTask} totalisent \textbf{zéro occurrence}. Un point gouverne l'ordonnancement :
\code{HopKind} est le goulot partagé -- conçu par l'Encadrement, mais consommé aussi par l'essaim
médiateur \emph{et} par la Colonne, de sorte que le premier incrément débloque les trois pistes.

\subsection{3.4 · Les nano-robots -- réels, mais sémantiquement vides}

La ligne de tenseur existe, la fusion est réelle, le pas de couplage est déterministe. Mais les priors
des templates valent \code{scale * sin(i)}, avec ce commentaire dans le source même :
\emph{« Not meaningful -- just distinguishable across templates »} (\code{genesis/registry.rs:81}). Il
n'y a \textbf{aucun poids de modèle de langage et aucun tokenizer} sur le chemin d'exécution. Ce n'est
pas un trou à combler, c'est \textbf{une phrase à ne jamais laisser tomber} : « nano-robots » désigne
des lignes de tenseur à priors vides, pas des agents.

Deux trous réels, en revanche. D'abord, \textbf{l'état initial est hors du domaine déclaré} :
l'initialisation et la remise à zéro placent le participant à l'origine, et le contrôle d'état ne
vérifie que la dimension et la finitude -- or avec la boîte de sûreté $[0{,}05\,;0{,}30]^{16}$,
\textbf{zéro est en dehors}. La contenance n'est rétablie qu'après la projection du premier pas, si bien
que l'affirmation du canari (« échantillonné sur la boîte ») ne décrit pas les états réellement
atteignables. \sz{XS--S} Ensuite, \textbf{la cardinalité déclarée est ignorée} : la genèse dérive le
nombre d'agents du graphe et jamais du plan, de sorte que \code{spec.population} n'est pas honoré sur la
branche dynamique -- au moins 128 agents éclosent là où le manifeste en demande quatre. C'est
précisément ce qui confond la mesure décrite ci-dessous. \sz{S}

\subsection{3.5 · L'essaim $\Phi$ -- le composant le plus honnête du dépôt}

Il faut le dire tel quel, car c'est rare : \textbf{la borne analytique $\beta_1 \approx 0{,}17$ est
explicitement réfutée dans le code lui-même}, à deux endroits, au profit de la seule borne
\emph{mesurée}. Mieux, \code{phi\_once} est un accesseur public réellement utilisé par le pas de
couplage -- \textbf{l'opérateur mesuré est donc bien l'opérateur appliqué}, ce qui est loin d'être
garanti dans ce genre de code. Le canari de non-expansivité fait 590 lignes, procède par bissection
adverse à la recherche de discontinuités, et son propre en-tête impose l'étiquette : \emph{« c'est un
FALSIFICATEUR, pas une preuve\dots{} il ne peut jamais certifier »}.

Le reste-à-faire est d'abord un défaut de concurrence, et il touche le chemin \emph{par défaut} :
l'assemblage à topologie fixe construit son dissolveur sur le \textbf{slot partagé} du contrôleur, alors
que le correctif -- le centroïde porté par l'ELS -- n'est appliqué que sur la branche dynamique, et
qu'il est toujours \textbf{non commité}. Or le gestionnaire de pas ne sérialise pas les appels par
participant : deux pas concurrents peuvent se rendre \textbf{le centroïde l'un de l'autre}, mode de
défaillance que la documentation du module décrit exactement. \sz{S--M} Viennent ensuite trois dettes de
mesure : \textbf{l'écart 0,366 contre 0,17 n'est pas réconcilié} -- la piste est nommée, mesurer les
étages un à un \sz{M} ; \textbf{la comparaison fixe-contre-dynamique est confondue}, la population et le
mélange de templates changeant en même temps que la fixité du routage, si bien que l'addendum
\textbf{retire lui-même son claim causal} \sz{S} ; et \textbf{le budget de canari par défaut n'a jamais
été exécuté}, le seul test le réduisant à la main -- on ignore s'il conclut ou s'il refuse de servir à
chaque démarrage. \sz{S}

\subsection{3.6 · Le méta-essaim -- l'écart le plus large entre le dit et le servi}

\textbf{Réfutation n\textsuperscript{o}1 : il sert encore le mock.} La vérification est empirique et
sans ambiguïté -- la liste des features par défaut est \code{["tls"]}, et l'arborescence de dépendances
du binaire de production ne trouve \textbf{aucune} trace d'efe-core. Une ligne de manifeste sans bloc
\code{controller} construit un participant simulé dont l'opérateur est littéralement
$\Phi(x) = \kappa \cdot x$. Les jalons M1.a à M1.d \emph{sont} codés, et bien codés -- mais derrière un
drapeau qui n'est ni activé par défaut, ni construit en intégration continue, ni rendu par la charte de
déploiement.

Le trou le plus grave n'est nommé par aucun document : \textbf{le contrôleur ne vérifie jamais à quel
moteur il parle}. Le service de capacités est servi, porte l'empreinte de configuration et le jeton
d'identité du moteur -- et n'a \textbf{aucun appelant}. Un manifeste déclarant des participants efe-core
peut donc être servi de bout en bout \textbf{par des mocks}, sans qu'aucune erreur ne se déclenche :
c'est exactement la précondition que le jalon C.3 de la Colonne suppose acquise. \sz{M} Dans le même
esprit, \textbf{le garde H1 servi ne peut pas échouer} : il existe, fail-closed, mais uniquement dans le
processus du participant, au démarrage, en mémoire, sous une feature que personne ne construit --
\emph{aucune évidence H1 ne traverse le fil}. La formulation honnête est donc : chaque participant réel
s'auto-contrôle au boot et refuse de servir ; \textbf{le contrôleur le croit sur parole}.

Les jalons déclarés qui restent sont massifs. \textbf{M2} (servi et reçu) est à \textbf{zéro ligne} : le
mot \emph{provenance} n'apparaît même pas dans le manifeste de dépendances du serveur, ce qui signifie
que \textbf{le méta-essaim n'écrit aucun reçu} -- la colonne vertébrale du projet s'arrête à sa porte.
\sz{L} \textbf{M3} (agrégateur robuste) est à zéro également : ni médiane, ni moyenne tronquée, ni Krum,
et l'opérateur uniforme est codé en dur côté serveur, la configuration n'étant que lue -- un participant
adverse pilote donc la moyenne. \sz{L} \textbf{M4} (certificat réel) est le manque porteur : le
certificat est un \code{Stub} déclaré valide, la crate \code{efe-prove} est citée dans les commentaires
mais \textbf{n'existe pas}, et le hachage de l'ensemble de sûreté retourne encore trente-deux octets
nuls. \sz{L} \textbf{M0} n'est pas fait non plus : la porte qui doit refuser un opérateur expansif à la
construction n'existe pas. \sz{XS}

S'y ajoutent quatre manques opérationnels. \code{require\_convergence} vaut \textbf{\code{false} par
défaut} et aucune règle ne le lie à un moteur réel : le point fixe parasite rendu comme « consensus »
reste atteignable avec un manifeste standard, alors qu'une seule règle de validateur suffirait. \sz{XS}
Le composant est \textbf{indéployable via la charte} livrée, qui rend inconditionnellement les
paramètres du mock et n'émet aucun bloc \code{controller}. \sz{M} Il serait de surcroît
\textbf{instable} : le canari de démarrage s'exécute \emph{avant} l'ouverture du port, à raison de
quelques secondes par évaluation, alors que la sonde de vivacité tue vers quarante secondes -- il
manque une sonde de démarrage. \sz{XS, mais bloquant} Enfin \textbf{aucun test ne pilote une ronde
complète contre de vrais participants} : le test de bout en bout ne conduit qu'\emph{un} pas, et le test
de convergence utilise \textbf{trois mocks}. \sz{M}

\subsection{3.7 · Le reçu, pile A -- l'actif-roi, et son défaut}

Il faut commencer par ce qui est solide, car c'est réel : environ 9\,600 lignes et 320 tests, une
chaîne HMAC par tenant, une chaîne N17 parallèle pour les appels de modèles, un pont entre les deux, et
une re-preuve hors ligne qui fonctionne. Le modèle de confiance est honnêtement borné dans ses propres
documents comme vérifiable par le détenteur de la clé, jamais comme sans confiance.

Mais trois défauts critiques doivent être fermés avant toute revendication. Le premier est le défaut de
troncature décrit en tête de ce document. Le deuxième est que \textbf{le résumé de requête n'est pas
scellé} alors que son commentaire de documentation le qualifie de « scellé » : c'est une simple colonne
JSON, et c'est pourtant elle que le vérificateur d'arité interroge -- si bien que qui écrit dans la base
réécrit l'arité déclarée et fait disparaître l'accusation. Le troisième est que la déclaration engagée
dans le n\oe{}ud de dispatch est \textbf{inouvrable} : l'empreinte engage la liste des fournisseurs,
mais celle-ci n'est persistée nulle part, de sorte que personne -- pas même le détenteur de la clé -- ne
peut reconstituer la préimage. La « déclaration inviolable » ne prouve donc rien à personne.

S'y ajoutent des manques de moindre gravité mais réels : les axes d'honnêteté (le barreau le plus
faible, les anomalies de cosignature et de modèle) sont \emph{calculés puis jetés} avant d'atteindre la
surface HTTP, alors que les invariants exigent qu'ils soient « toujours rapportés, jamais silencieux » ;
le registre de n\oe{}uds n'étant jamais câblé, \textbf{toute} revendication R3 se dégrade
systématiquement en R4 ; la clé de développement, publique, est utilisée lorsque la frontière est
désactivée ; aucune version de clé n'est enregistrée, si bien qu'une rotation rendrait tous les reçus
antérieurs indistinguables d'une falsification ; et le champ \code{audit\_chain\_hash}, servi dans
chaque réponse de mission et documenté comme une « preuve d'intégrité BLAKE3 », est en réalité
l'identifiant de la mission rendu en hexadécimal.

\begin{these}
\textbf{La bonne nouvelle, et elle est structurelle.} La vérifiabilité par un tiers est bien plus proche
qu'il n'y paraît. Les empreintes de contenu -- celle de chaque enregistrement comme celle de chaque
feuille -- sont \textbf{sans clé} : n'importe qui peut déjà les recalculer depuis les enregistrements
persistés. Seuls le HMAC de chaque maillon et le pli de racine exigent le secret. Il suffit donc de
\textbf{signer la racine} avec une paire de clés asymétrique pour qu'un tiers vérifie sans rien
détenir -- et la bibliothèque ed25519 est \textbf{déjà présente} dans la crate, avec toute la machinerie
de cosignature. La cryptographie est une affaire de taille S ; la garde et la rotation des clés, de
taille M. C'est le meilleur rapport valeur-effort de tout le dépôt.
\end{these}

\subsection{3.8 · Le reçu, pile B -- le moat, complet mais sans producteur}

\textbf{Réfutation n\textsuperscript{o}2 : les sept tranches sont construites.} Ma note antérieure
indiquait que la console, l'export forensique et le moteur d'évidence restaient à faire. C'est faux :
les neuf crates existent, totalisent environ 11\,700 lignes et \textbf{245 tests}. Mieux, la
revendication centrale est \emph{tenue et testée} : le vérificateur minimal a une arborescence de
dépendances littéralement vide, et un test de bout en bout recompile réellement la source vendorée hors
ligne puis re-prouve un dossier. « Un tiers vérifie sans nous faire confiance » est ici un test qui
passe, pas une diapositive.

Le reste-à-faire est donc d'une autre nature : il ne s'agit plus d'écrire de la cryptographie, mais de
brancher ce qui existe. Et le trou est béant : \textbf{il n'y a aucun producteur}. Les « prises » que la
feuille de route inventoriait comme un proxy de passerelle, un crochet d'orchestration et un
intercepteur d'outils sont en réalité \textbf{quatre structures de conversion}. Rien, dans tout le
dépôt, ne construit un graphe de provenance en dehors des exemples et des tests. S'y ajoutent trois
manques opérationnels sérieux : la décision humaine scellée par un opérateur est \textbf{perdue au
redémarrage}, parce que la console n'appelle jamais la fonction de persistance ; \textbf{aucun ancrage
n'est jamais produit}, de sorte que la racine que l'auditeur doit comparer \emph{voyage à l'intérieur du
dossier qu'on lui tend} ; et la seule route qui scelle dans le graphe n'a aucune protection contre la
soumission inter-site, ce qui permettrait de sceller une décision humaine cryptographiquement
indistinguable d'une vraie. L'authentification, enfin, se réduit à un jeton de démonstration en
mémoire. Le seul jalon formellement restant, \textbf{S7}, n'est pas du code : c'est une porte humaine,
un essai d'utilisabilité et la conversion d'un partenaire de conception.

\subsection{3.9 · La Couche 2 servie -- complète, mais pas sûre en production}

Les cinq jalons M0 à M4 ont du code réel, des tests réels et des documents réels : le test-gardien qui
fait échouer la vérification quand un appel a été omis, la triade de démarrage, le décompte des coûts en
microcentimes, la clé Gemini déplacée dans un en-tête, la fermeture du contournement SSRF sur les deux
copies du garde, le document de modèle de confiance et un manuel d'opérateur de 164 lignes existent
tous. Rien de tout cela n'est servi par défaut, ce qui est un choix assumé.

Ce qui reste tient en deux défauts de sûreté et un déficit de preuve. Le premier défaut est que
\textbf{la rédaction des données personnelles est fail-open à la compilation} : elle est conditionnée à
une feature qui n'a \emph{aucune contrepartie} et que la porte de démarrage ne vérifie pas. Un binaire
construit avec exactement la triade que cette porte exige passe donc le contrôle, affiche « préflight
frontière OK », et expédie tous les prompts en clair sans le moindre signal. Le second est que
\textbf{le plafond de budget est inopérant} pour tout tenant dont la politique ne porte pas la mention
de plafond : une requête peut alors réclamer un budget de l'ordre de quarante-deux millions de dollars,
ce qui est précisément le cas que le critère d'acceptation du jalon M4 demandait de fermer -- et il
n'existe aucun test de cette fonction. Le déficit de preuve, lui, est que le test de bout en bout ne
couvre \textbf{qu'une mission sur cinq} : ni le relais ni la cascade n'ont de test HTTP. S'y ajoutent
l'absence de limitation de débit, l'absence de quota cumulé, l'absence totale de journalisation de
requête sur le seul chemin qui dépense de l'argent réel, et une charte de déploiement dont la politique
réseau bloque en réalité les fournisseurs qu'elle prétend appeler.

\subsection{3.10 · Le runtime nano -- une bibliothèque que rien n'exécute}

Les sept défauts relevés lors de l'audit précédent sont \textbf{tous confirmés et aucun n'est corrigé} ;
trois se révèlent plus graves que décrits. Le plus sérieux concerne la canonicalisation JSON : les
\emph{clés} d'objet sont écrites \textbf{sans échappement} alors que les valeurs, elles, sont échappées.
On peut donc forger deux entrées distinctes produisant la même empreinte. Or cette fonction n'alimente
pas seulement le cache, comme je l'avais d'abord estimé : elle alimente l'empreinte de charge utile de
la \textbf{chaîne d'audit} elle-même. Ce n'est donc pas une nuisance de cache, c'est une
\textbf{primitive de forge de provenance qui se re-vérifie verte}. Le deuxième défaut aggravé est une
\textbf{dérive entre canonicaliseurs} : la pile de provenance échappe correctement les clés, le runtime
nano non, et un commentaire affirme pourtant que les deux sont identiques -- exactement la classe de bug
de la duplication du garde de souveraineté déjà survenue. Le troisième est que les écritures
concurrentes d'audit pour un même tenant relâchent le verrou avant l'ajout, si bien que des traces sont
\textbf{silencieusement perdues} -- et l'exécuteur, lui, est parallèle par tenant, ce qui rend le cas
non pas théorique mais systématique. Deux commentaires de documentation affirment le contraire, et le
cas dangereux n'est pas testé.

À cela s'ajoute une composition redoutable : l'écriture d'audit est fail-open et la source de clé
retourne une erreur quand la variable d'environnement maîtresse est absente. Déployer sans cette
variable revient donc à \textbf{désactiver entièrement l'audit}, sans aucun signal, le processus
paraissant en parfaite santé. Il n'existe aucune porte de démarrage. Le poids d'équité de
l'ordonnanceur reste inerte -- et son test d'acceptation déclaré n'a jamais été écrit --, les
comparaisons de vérification ne sont pas à temps constant, l'empreinte de pierre tombale est calculée
mais \textbf{jamais vérifiée}, et l'interface d'effacement n'a aucun appelant en production.

\subsection{3.11 · Le banc de mesure -- des statistiques vraies, des mesures non régressibles}

\textbf{Réfutation n\textsuperscript{o}3} : il n'existe \textbf{aucun} test live marqué
\code{\#[ignore]} dans les crates de banc. Les pilotes live sont des exemples lancés \emph{à la main}.
Le c\oe{}ur analytique, lui, est solide et je l'ai vérifié : 118 tests verts, non-infériorité appariée,
bootstrap déterministe, correction de Holm, constantes gelées avant l'expérience. Ses verdicts
\emph{négatifs} -- l'agrégation ne comble pas l'écart, aucune des trois cellules de terrain ne creuse --
sont les artefacts les plus crédibles du dépôt.

Le problème n'est pas la statistique, c'est l'infrastructure autour. Le seul test de bout en bout du
harnais est \textbf{compilé hors intégration continue} par une condition de feature, ce qui se corrige
en une ligne. \textbf{Aucun nombre publié n'est reproductible depuis un script commité} : ni dans la
CI, ni dans un fichier de commandes, ni dans le document de benchmarks. Et la porte de test gold, censée
protéger la qualité, garde en réalité un simulateur d'écho avec un fournisseur d'embeddings factice,
tandis que son balayage de dérive retourne « à jour » pour tout agent quel que soit son état. Le levier
décisif reste les \textbf{cassettes d'enregistrement et de rejeu} : elles transforment chaque chiffre
des six rapports d'une assertion de fichier markdown en un résultat que la CI re-dérive.

\section{4 · Les cinq fils transverses}

\textbf{Le fossé construit-servi.} Le méta-essaim, la Couche 2, le runtime nano et la pile B sont
écrits ; presque rien ne tourne. Le travail restant relève plus souvent du \emph{câblage} que de la
construction, ce qui est une bonne nouvelle sur l'effort et une mauvaise sur ce qu'on peut revendiquer
aujourd'hui.

\textbf{Le motif du « vert malhonnête ».} Cinq occurrences indépendantes ont été trouvées : la
vérification sans comparaison de racine, le service de capacités jamais appelé, le balayage de dérive
toujours favorable, le certificat déclaré valide par construction, et l'audit fail-open. C'est la
\emph{classe de bug signature} du dépôt -- celle qu'il faut désormais chercher systématiquement, parce
qu'elle ne casse rien : elle rassure à tort.

\textbf{L'ancrage externe manque aux deux piles.} Ni la pile A ni la pile B ne peut détecter une
troncature de queue par le détenteur de la clé. C'est le même trou, deux fois, et c'est celui qui sépare
« re-vérifiable » de « opposable ».

\textbf{Rien de critique n'est en intégration continue} : ni le vrai moteur latent, ni le banc, ni les
tests de base de données durable. Les revendications reposent donc sur des exécutions manuelles.

\textbf{Les canonicaliseurs dérivent} -- avec un commentaire qui affirme l'inverse. Deux implémentations
qui doivent s'accorder au bit près et qui ne s'accordent plus.

\section{5 · L'ordre recommandé}

\begin{center}\small
\begin{tabular}{@{}r>{\raggedright\arraybackslash}p{9.5cm}c>{\raggedright\arraybackslash}p{3.9cm}@{}}
\toprule
& \textbf{Action} & \textbf{Taille} & \textbf{Ce qu'elle ferme} \\
\midrule
1 & Comparer la racine dans \code{verify} & \sz{S} & le reçu ment aujourd'hui \\
2 & Sceller le résumé de requête, persister la déclaration & \sz{S} & l'arité redevient une preuve \\
3 & Échapper les clés dans la canonicalisation & \sz{XS} & la forge verte \\
4 & Porte de démarrage : PII dans la triade, clé nano exigée & \sz{S} & deux fail-open \\
5 & Plafond de budget serveur par défaut, avec test & \sz{S} & le critère M4 non tenu \\
6 & Rendre le filtrage de moteurs exhaustif & \sz{XS} & avant toute extension \\
\midrule
7 & Signature asymétrique sur la racine & \sz{S+M} & la vérifiabilité \textbf{tierce} \\
8 & Cassettes d'enregistrement et de rejeu & \sz{L} & les chiffres deviennent régressibles \\
9 & L'Encadrement (C.1.a$\rightarrow$d) & \sz{L} & la première capacité neuve \\
\bottomrule
\end{tabular}
\end{center}

Les six premiers items totalisent \textbf{moins d'une semaine} et font passer le reçu de « il se
re-prouve » à « il prouve ce qu'il dit ». Tant qu'ils ne sont pas faits, chaque brique supplémentaire
s'appuie sur une fondation qui renvoie \code{verified: true} à tort. Le septième est le seul qui change
la \emph{nature} de l'actif : il le rend opposable à un tiers, ce qui est précisément ce que le marché
de la conformité achète.

\enlargethispage{3.2\baselineskip}
\section{6 · Ce qu'il ne faut jamais revendiquer, en l'état}

Que le reçu détecte une troncature -- il ne la détecte pas aujourd'hui. Qu'il soit sans confiance ou
vérifiable par un tiers -- il est tamper-evident pour le détenteur de la clé, sans ancrage externe. Que
le méta-essaim coordonne de vrais essaims en production -- il sert le mock par défaut. Qu'il coordonne
des petits modèles de langage -- les priors sont un simple sinus, auto-qualifiés de « non signifiants »,
et aucun poids de modèle ni tokenizer ne se trouve sur le chemin. Que la ronde servie vérifie
l'hypothèse de non-expansivité -- le contrôleur est construit en mode indulgent. Que le runtime nano
soit servi, que son ordonnanceur assure l'équité, ou que l'effacement au titre du RGPD fonctionne de
bout en bout. Que les chiffres publiés soient re-prouvés cryptographiquement -- ce sont des comptages de
drapeaux. Et que l'agrégation apporte quoi que ce soit : elle a été \textbf{réfutée par les benchmarks
du projet lui-même}, et c'est un actif d'honnêteté, pas une faiblesse à masquer.

\vspace{4pt}
{\color{grey}\hrule height 0.4pt}
\vspace{2pt}
{\footnotesize\color{grey}Analyse produite par relecture du code au 3 août 2026, non des documents de
conception. Les tailles (XS à XL) sont des estimations d'effort, pas des engagements. Trois
affirmations antérieures de l'auteur y sont réfutées et signalées comme telles.}

\end{document}
